% Options for packages loaded elsewhere
\PassOptionsToPackage{unicode}{hyperref}
\PassOptionsToPackage{hyphens}{url}
\PassOptionsToPackage{dvipsnames,svgnames,x11names}{xcolor}
\documentclass[
  11pt,
]{ctexart}
\usepackage{xcolor}
\usepackage[margin=2.5cm]{geometry}
\usepackage{amsmath,amssymb}
\setcounter{secnumdepth}{-\maxdimen} % remove section numbering
\usepackage{iftex}
\ifPDFTeX
  \usepackage[T1]{fontenc}
  \usepackage[utf8]{inputenc}
  \usepackage{textcomp} % provide euro and other symbols
\else % if luatex or xetex
  \usepackage{unicode-math} % this also loads fontspec
  \defaultfontfeatures{Scale=MatchLowercase}
  \defaultfontfeatures[\rmfamily]{Ligatures=TeX,Scale=1}
\fi
\usepackage{lmodern}
\ifPDFTeX\else
  % xetex/luatex font selection
\fi
% Use upquote if available, for straight quotes in verbatim environments
\IfFileExists{upquote.sty}{\usepackage{upquote}}{}
\IfFileExists{microtype.sty}{% use microtype if available
  \usepackage[]{microtype}
  \UseMicrotypeSet[protrusion]{basicmath} % disable protrusion for tt fonts
}{}
\makeatletter
\@ifundefined{KOMAClassName}{% if non-KOMA class
  \IfFileExists{parskip.sty}{%
    \usepackage{parskip}
  }{% else
    \setlength{\parindent}{0pt}
    \setlength{\parskip}{6pt plus 2pt minus 1pt}}
}{% if KOMA class
  \KOMAoptions{parskip=half}}
\makeatother
\setlength{\emergencystretch}{3em} % prevent overfull lines
\providecommand{\tightlist}{%
  \setlength{\itemsep}{0pt}\setlength{\parskip}{0pt}}
\newcommand{\blank}{\underline{\hspace{2.8em}}}
\usepackage{bookmark}
\IfFileExists{xurl.sty}{\usepackage{xurl}}{} % add URL line breaks if available
\urlstyle{same}
\hypersetup{
  colorlinks=true,
  linkcolor={blue},
  filecolor={Maroon},
  citecolor={Blue},
  urlcolor={blue},
  pdfcreator={LaTeX via pandoc}}

\author{}
\date{}

\begin{document}

\section{《虚拟地理环境》期末考试}\label{ux865aux62dfux5730ux7406ux73afux5883ux671fux672bux8003ux8bd5}

\begin{itemize}
\tightlist
\item
  福建师范大学地理科学学院
\item
  课程：《虚拟地理环境》
\item
  任课教师：江辉仙
\item
  考试时间：2026年6月3日 8:20 共120分钟
\item
  考试专业：2023级 地理信息科学专业 专业选修课
\item
  考试形式：线下 闭卷
\end{itemize}

题目内容为回忆版，可能与实际试题有所出入。已核对试题顺序及具体考察内容无误。\href{https://github.com/Xuuyuan}{@Xuuyuan}

\subsection{题目}\label{ux9898ux76ee}

\subsubsection{填空题（10题，每空1分，共20分）}\label{ux586bux7a7aux989810ux9898ux6bcfux7a7a1ux5206ux517120ux5206}

\begin{enumerate}
\def\labelenumi{\arabic{enumi}.}
\tightlist
\item
  交互设计的两项核心内容是\blank{}和\blank{}。
\item
  如果在一个文件中定义了一个节点名称，可以在相同的文件中使用\blank{}来重用已经定义的节点名称。原来被定义名称的节点称做\blank{}，每次被重用的节点则叫做\blank{}。
\item
  VRML语言中以\blank{}开头的内容为注释。
\item
  ROUTE用于将一个节点的\blank{}事件传递给另一个节点的\blank{}事件。
\item
  三维几何建模的常见模型有\blank{}、\blank{}、\blank{}。
\item
  VRML中的字段可以分成两类，一类是字段只包含一个单独值，开头的名称为\blank{};另一类是包含多重值，开头名称为\blank{}。
\item
  VRML是\blank{}类型语言，因此它可以在记事本等程序中打开。
\item
  虚拟地理环境的三个发展阶段：\blank{}、\blank{}、\blank{}。
\item
  由于观察位置不同而导致同一物体在视网膜或图像中的位置差异称为\blank{}。
\item
  基于物理设备跟踪的输入设备有道具、\blank{}、\blank{}。
\end{enumerate}

\subsubsection{名词解释（4题，每题4分）}\label{ux540dux8bcdux89e3ux91ca4ux9898ux6bcfux98984ux5206}

\begin{enumerate}
\def\labelenumi{\arabic{enumi}.}
\tightlist
\item
  虚拟地理环境
\item
  三维建模技术
\item
  地理知识图
\item
  数字地球
\end{enumerate}

\subsubsection{简答题（4题，每题8分）}\label{ux7b80ux7b54ux98984ux9898ux6bcfux98988ux5206}

\begin{enumerate}
\def\labelenumi{\arabic{enumi}.}
\tightlist
\item
  简述3DGIS与VRGIS的区别。
\item
  简述地理信息可视化的含义、表现形式和过程。
\item
  简述虚拟地理环境的层面结构。
\item
  简述在VRML中管理大文件的方法（所用的节点和使用方法）。
\end{enumerate}

\subsubsection{论述题（2题，每题16分）}\label{ux8bbaux8ff0ux98982ux9898ux6bcfux989816ux5206}

\begin{enumerate}
\def\labelenumi{\arabic{enumi}.}
\tightlist
\item
  请说明虚拟现实、增强现实、混合现实各自的特点和区别。
\item
  根据你或你们小组的系统设计，说明虚拟地理系统的目的、特色及设计过程。
\end{enumerate}

\end{document}
